% Complete reconstruction of the local tame-p transcription audited by
% root and D, with the equivariant Matlis calculation retained explicitly.
\begin{lemma}[Proved: A uniform ordinary Frobenius defect in the Kummer rings]
\label{D:Kummer-uniform-Frobenius-defect}
In the local notation of Theorem~\ref{D:local-tame-p-linear-H1}, one
fixed $\mathfrak m$-primary ideal $J\subset R$ has the following
property for every $n$ prime to $p$. For every maximal Cohen--Macaulay
$R_n$-module $E$, multiplication by every element of $JR_n$ factors
through a finite free module. The cokernel of
$F_{R_n}^*E\to F_{R_n}^{[*]}E$ is annihilated by the ideal generated
by the $p$th powers of the elements of $J$.
\end{lemma}
\begin{proof}
Use the fixed ideal $J_0$ of base Jacobian minors displayed in the
proof below. Its extension belongs to the Jacobian ideal of every
$R_n$, independently of $n$. In this equicharacteristic complete
setting the Jacobian ideal is the sum of the Kahler differents for
regular Noether normalizations, and each Kahler different belongs to
the corresponding Noether different
\citep[Section 2, pp.~2--3, and Section 3(3.1), Examples 3.2--3.3,
pp.~6--7]{C_IyengarTakahashiStable}.
We recall the relevant factorization. For a finite regular
normalization $A\subset R_n$, write
$\eta=\sum u_i\otimes v_i$ in the annihilator of the kernel of
$R_n\otimes_A R_n\to R_n$, and let $z=\sum u_iv_i$.
The maps
\[
 E\longrightarrow R_n\otimes_A E\longrightarrow E,
 \qquad e\longmapsto\sum u_i\otimes v_ie,
 \qquad b\otimes e\longmapsto be
\]
are $R_n$-linear and compose to $z\operatorname{id}_E$.
The middle module is free over $R_n$, because $E$ is maximal
Cohen--Macaulay and therefore free over the regular local ring $A$.
Taking finite sums gives the assertion for $J_0R_n$.
Frobenius pullback of the factorization makes $z^p$ times the identity
on $F^*E$ factor through a free module. Taking reflexive hulls of the
maps, the free middle module remains unchanged, and shows that
$z^pF^{[*]}E$ lands in the image of $F^*E$. We may take $J=J_0$,
or any fixed $\mathfrak m$-primary subideal. This proves both claims.
\end{proof}

\begin{theorem}[Proved: A uniform local tame bound in coefficient characteristic $p$]
\label{D:local-tame-p-linear-H1}
Let $(R,\mathfrak m,k)$ be a complete normal local surface ring over an
algebraically closed field of characteristic $p>0$. Let $D$ be a fixed
reduced Weil divisor, and let $Q_D$ be the quotient of the fundamental
group of $\operatorname{Spec}R\setminus(D\cup\{\mathfrak m\})$
which requires tameness at the original height-one components of $D$.
There is a constant $c_{R,D}$ such that every finite-dimensional module
$M$ over a finite field of characteristic $p$ satisfies
\[
 \dim H^1(Q_D,M)\le c_{R,D}\dim M.
\]
The same bound applies to the image in local $H^1$ of classes from any
global fixed-original-divisor tame quotient. The assertion also holds
for excellent strictly henselian surface rings by completion.
\end{theorem}
\begin{proof}
We first use $\mathbf F_p$ coefficients. Choose a fixed reduced principal
divisor $\operatorname{div}(a)$ containing $D$. A general element of
the ideal of $D$ has this property: at each prescribed height-one prime
we choose valuation one, and on the regular puncture the remaining
divisor is generically reduced. We may enlarge $D$ by these auxiliary
components temporarily. Restriction to their complement is injective
on $H^1$, because a connected cover of a normal integral scheme remains
connected on a dense open.

For every $n$ prime to $p$, put
\[
 R_n=R[s]/(s^n-a),\qquad G=\mu_n.
\]
This is a normal Cohen--Macaulay local domain. Indeed it is finite free
over $R$ and therefore Cohen--Macaulay; at a component of $(a)$ it is
an Eisenstein extension of a regular DVR, and elsewhere on the puncture
it is etale. Serre's criterion proves normality, and reduction modulo
the maximal ideal proves locality. These arguments also show that its
puncture is regular.

We require one annihilator independent of $n$. Present
$R=k[[x_1,\ldots,x_v]]/(f_1,\ldots,f_t)$, of codimension $c=v-2$.
Let $I_j$ denote the ideal of $j$-minors. The fixed ideal
\[
 J_0=I_{c+1}\bigl(\operatorname{Jac}(f_1,\ldots,f_t,a)\bigr)
       +aI_c\bigl(\operatorname{Jac}(f_1,\ldots,f_t)\bigr)
       \subset R
\]
is $\mathfrak m$-primary or the unit ideal. Outside the closed point
the ring is regular, and the reduced curve $(a)$ is regular at its
generic points. The Jacobian matrix for $R_n$ has final row
$(-\partial a/\partial x_i,ns^{n-1})$. Its minors contain $J_0R_n$:
the first summand uses no $s$-column, and the second is contained in
$s^{n-1}I_c(\operatorname{Jac}f)$ because $a=s^n$ and $n$ is a unit.
For a finite generically separable regular Noether normalization, the
relative Jacobian elements are completely stable big test elements by
the Lipman--Sathaye theorem in
\citet[printed p.~207]{C_Hochster2007}. The complete-local Jacobian is
the sum of these relative Kahler differents, as in
\citet[Section 3(3.1), Examples 3.2--3.3]{C_IyengarTakahashiStable}.
This proves the assertion for our complete rings; the affine version
is also stated on printed p.~213 of Hochster's notes. Thus the same
fixed $\mathfrak m$-primary ideal $J\subset R$ annihilates the tight
closure of zero in the injective hull over every $R_n$. We may take
$J=J_0$, or any fixed $\mathfrak m$-primary subideal.

Choose $n$ divisible by the orders of the finitely many original tame
inertia images in the representation $M$. The Kummer pullback extends
to a local system $M_n$ on the puncture of $R_n$, by the DVR tame
calculation and purity on that regular puncture. The same $n$ works for
every class in $H^1(Q_D,M)$. Such a class is represented by affine
monodromy in $M\rtimes\operatorname{im}(Q_D\to\operatorname{GL}(M))$.
Its tame inertia is a group of order prime to $p$ and has trivial
intersection with the additive $p$-group $M$. Its projection therefore
has the same order, already dividing $n$. This explains why the cover
used below is independent of the cohomology class.

Write $E_n$ for the unit reflexive extension of
$\mathcal O\otimes M_n$, and set $H_n=H^2_{\mathfrak m_n}(E_n)$.
All constructions are $G$-equivariant. The proof of
Theorem~\ref{D:local-p-linear-H1} gives a Frobenius-injective quotient
$H_n/H_{n,\rm nil}$ with $JH_{n,\rm nil}=0$. We take Matlis duals
over the fixed ring $R$, rather than estimate ranks over $R_n$. Put
\[
 P=\operatorname{Hom}_R(E_n,\omega_R),\qquad
 N=D_R(H_n/H_{n,\rm nil}).
\]
Then $JP\subset N\subset P$. Since $|G|$ is invertible, invariants
are exact. The invariant part
\[
 P^G=\operatorname{Hom}_R(E_n^G,\omega_R)
\]
is maximal Cohen--Macaulay of rank $r$ over $R$: $E_n^G$ has rank $r$
by generic Galois descent. We have
$JP^G\subset N^G\subset P^G$. A fixed parameter ideal
$Q\subset\mathfrak mJ$ therefore gives
\[
 \mu_R(N^G)\le\operatorname{length}(P^G/QP^G)
 =r\operatorname{length}(R/Q).
\]
The Frobenius-fixed invariant vectors inject into the invariant socle
of $H_n/H_{n,\rm nil}$. Exact invariants and Matlis duality identify
its dimension with $\mu_R(N^G)$. Thus the invariant part of
$H^1$ on the upstairs puncture has dimension at most $c_{R,D}r$,
with no factor $n$.

Restriction to the Kummer cover is injective on $H^1$ in characteristic
$p$, by restriction--corestriction, since its degree $n$ is invertible.
Its image is invariant, and the preceding estimate applies to every
class under consideration. This proves the theorem over $\mathbf F_p$.
Restriction of scalars proves the assertion for every finite field of
characteristic $p$.

For a strictly henselian ring we use the completion surjection on
fundamental groups proved in Theorem~\ref{D:local-p-linear-H1}.
The original height-one tame condition remains tame after completion,
so the same injection and estimate apply. The image version for global
classes follows by exactly the same affine-monodromy argument; it does
not assert a bound for every unrestricted local class at a divisor.
\end{proof}
